\documentclass[11pt, a4paper]{article}

% bfw-style.tex — gemeinsame Style-Definitionen für alle Handouts/Aufgabenhefte
% Voraussetzung: Kompilation mit xelatex (wegen fontspec)

\usepackage[ngerman]{babel}
\usepackage{geometry}
\geometry{a4paper, top=2.2cm, bottom=2.2cm, left=2.3cm, right=2.3cm}

\usepackage{fontspec}
% Fallback-Kette: Source Sans Pro -> Calibri -> Segoe UI -> Latin Modern Sans
% (Latin Modern ist immer mit MiKTeX vorhanden)
\IfFontExistsTF{Source Sans Pro}{%
  \setmainfont{Source Sans Pro}[Ligatures=TeX]%
}{\IfFontExistsTF{Calibri}{%
  \setmainfont{Calibri}[Ligatures=TeX]%
}{\IfFontExistsTF{Segoe UI}{%
  \setmainfont{Segoe UI}[Ligatures=TeX]%
}{%
  \setmainfont{Latin Modern Sans}[Ligatures=TeX]%
}}}
\IfFontExistsTF{Source Code Pro}{%
  \setmonofont{Source Code Pro}[Scale=0.92]%
}{\IfFontExistsTF{Consolas}{%
  \setmonofont{Consolas}[Scale=0.92]%
}{%
  \setmonofont{Latin Modern Mono}[Scale=0.92]%
}}

\usepackage{xcolor}
% Farbpalette: hell, freundlich, modern
\definecolor{bfwPrimary}{HTML}{0EA5A4}    % Petrol/Türkis - Primär
\definecolor{bfwPrimaryDark}{HTML}{0F766E} % dunkler Petrol für Headings
\definecolor{bfwAccent}{HTML}{F59E0B}     % Amber - Akzent
\definecolor{bfwSuccess}{HTML}{10B981}    % Grün - Erfolg / Lösung
\definecolor{bfwWarn}{HTML}{F97316}       % Orange - Achtung
\definecolor{bfwInk}{HTML}{0F172A}        % fast Schwarz
\definecolor{bfwInkSoft}{HTML}{475569}    % gedämpftes Grau
\definecolor{bfwBgSoft}{HTML}{F8FAFC}     % off-white
\definecolor{bfwBgTeal}{HTML}{ECFEFF}     % helles Türkis
\definecolor{bfwBgAmber}{HTML}{FEF3C7}    % helles Gelb
\definecolor{bfwBgRose}{HTML}{FFE4E6}     % helles Rosé für Eselsbrücken

\usepackage[colorlinks=true, linkcolor=bfwPrimaryDark, urlcolor=bfwPrimaryDark]{hyperref}
\usepackage{enumitem}
\setlist[itemize]{leftmargin=*, itemsep=2pt, topsep=2pt}
\setlist[enumerate]{leftmargin=*, itemsep=2pt, topsep=2pt}

\usepackage[most]{tcolorbox}
\usepackage{booktabs}
\usepackage{tikz}
\usetikzlibrary{decorations.pathmorphing, positioning, shapes.geometric, arrows.meta}

% Merkkästchen — wichtige Definition / Kernpunkt
\newtcolorbox{merksatz}[1][]{%
  enhanced, breakable,
  colback=bfwBgTeal,
  colframe=bfwPrimary,
  boxrule=0.6pt,
  arc=4pt,
  left=10pt, right=10pt, top=8pt, bottom=8pt,
  fonttitle=\bfseries\color{bfwPrimaryDark},
  title={\faLightbulb\ Merksatz},
  #1
}

% Eselsbrücke — Mnemoniks
\newtcolorbox{eselsbruecke}[1][]{%
  enhanced, breakable,
  colback=bfwBgRose,
  colframe=bfwAccent,
  boxrule=0.6pt,
  arc=4pt,
  left=10pt, right=10pt, top=8pt, bottom=8pt,
  fonttitle=\bfseries\color{bfwWarn},
  title={Eselsbrücke},
  #1
}

% Beispiel-Box
\newtcolorbox{beispiel}[1][]{%
  enhanced, breakable,
  colback=bfwBgSoft,
  colframe=bfwInkSoft,
  boxrule=0.4pt,
  arc=3pt,
  left=10pt, right=10pt, top=8pt, bottom=8pt,
  fonttitle=\bfseries\color{bfwInk},
  title={Beispiel},
  #1
}

% Lösungs-Box (für Aufgabenhefte)
\newtcolorbox{loesung}[1][]{%
  enhanced, breakable,
  colback=bfwBgAmber!40,
  colframe=bfwSuccess,
  boxrule=0.6pt,
  arc=3pt,
  left=10pt, right=10pt, top=8pt, bottom=8pt,
  fonttitle=\bfseries\color{bfwSuccess!70!black},
  title={Lösung},
  #1
}

% Glossar-Eintrag
\newcommand{\glossareintrag}[2]{%
  \par\noindent\textbf{\color{bfwPrimaryDark}#1} \enspace #2\par\smallskip
}

% Section-Stil — etwas farbiger als Standard
\usepackage{titlesec}
\titleformat{\section}
  {\Large\bfseries\color{bfwPrimaryDark}}
  {\thesection}{0.6em}{}
\titleformat{\subsection}
  {\large\bfseries\color{bfwPrimary}}
  {\thesubsection}{0.5em}{}
\titleformat{\subsubsection}
  {\normalsize\bfseries\color{bfwInk}}
  {\thesubsubsection}{0.4em}{}

% TikZ-Stil "skizzenhaft" — etwas weicher als Rough.js, aber stabil
\tikzset{
  roughbox/.style = {
    draw=bfwPrimaryDark, thick, fill=bfwBgTeal,
    rounded corners=4pt, inner sep=6pt
  },
  roughaccent/.style = {
    draw=bfwAccent, thick, fill=bfwBgAmber,
    rounded corners=4pt, inner sep=6pt
  },
  roughline/.style = {
    draw=bfwInkSoft, thick, -{Stealth[length=2.5mm]}
  }
}

% Bullet-Symbole (bewusst kein FontAwesome — vermeidet Font-Installations-Probleme)
\newcommand{\faLightbulb}{$\bullet$}
\newcommand{\faBookmark}{$\bullet$}

% Header/Footer
\usepackage{fancyhdr}
\pagestyle{fancy}
\fancyhf{}
\renewcommand{\headrulewidth}{0pt}
\fancyfoot[L]{\small\color{bfwInkSoft}\thepage}
\fancyfoot[R]{\small\color{bfwInkSoft} BFW M-064 Beschaffung im Büromanagement}


\title{\vspace*{-1.5cm}\Huge\bfseries\color{bfwPrimaryDark}Tag~8: Bezahlung von Rechnungen\\[0.4em]
  \large\color{bfwInkSoft}\textnormal{Geld, Zahlungsarten, SEPA, Karten, Onlinebanking, E-Commerce, Skonto}}
\author{\textsf{Beschaffung im Büromanagement}\\
\textsf{\small\copyright\ Can Siebert\,\textperiodcentered\,2026}}
\date{21.05.2026}

\begin{document}
\maketitle
\thispagestyle{empty}

\vspace{1em}
\begin{merksatz}[title={\bfseries Worum geht es heute?}]
Die Rechnung ist geprüft --- jetzt geht es ans Zahlen. Heute lernen wir die wichtigsten
Zahlungsarten kennen: bar, halbbar, bargeldlos. Schwerpunkt sind SEPA-Verfahren,
Zahlungskarten, Onlinebanking und die typischen Bezahlmethoden im E-Commerce.
Außerdem rechnen wir Skonto --- die wichtigste kaufmännische Berechnung der
Beschaffung und ein verlässlicher Punktelieferant in der IHK-Klausur Teil~2.
\end{merksatz}

\tableofcontents
\vspace{1em}
\hrule
\vspace{1em}

\section{Geldfunktionen und Geldarten}

\subsection{Drei Funktionen des Geldes}

\begin{itemize}[itemsep=2pt]
  \item \textbf{Tauschmittel}: Geld erleichtert den Tausch (statt Naturaltausch).
  \item \textbf{Wertaufbewahrung}: Geld kann gespart und für späteren Konsum aufbewahrt werden.
  \item \textbf{Wertmaßstab / Recheneinheit}: Geld misst den Wert von Gütern und ermöglicht Preisvergleiche.
\end{itemize}

\subsection{Geldarten}

\begin{itemize}[itemsep=2pt]
  \item \textbf{Bargeld}: Münzen und Banknoten --- gesetzliches Zahlungsmittel.
  \item \textbf{Buchgeld (Giralgeld)}: Guthaben auf Konten --- wird durch Überweisung, Lastschrift, Karte transferiert.
  \item \textbf{Elektronisches Geld}: digitale Werteinheiten (z.\,B.\ Prepaid-Karten, einige Wallet-Lösungen).
\end{itemize}

\section{Zahlungsarten}

\begin{center}
\begin{tabular}{p{3.5cm} p{10.5cm}}
\toprule
\textbf{Art} & \textbf{Beispiel}\\
\midrule
Bar & Kassenzahlung mit Münzen/Scheinen\\
Halbbar & Bareinzahlung auf ein Bankkonto, Postanweisung\\
Bargeldlos & SEPA-Überweisung, SEPA-Lastschrift, Kartenzahlung, PayPal\\
\bottomrule
\end{tabular}
\end{center}

\section{Bargeldlose Zahlungsformen}

\subsection{SEPA-Überweisung}

Der \emph{SEPA-Raum} (Single Euro Payments Area) umfasst 36~Länder. Jede Überweisung
nutzt \textbf{IBAN} (International Bank Account Number) --- in Deutschland 22~Stellen
(z.\,B.\ DE89~3704~0044~0532~0130~00). Standard-Überweisung dauert maximal einen Bankarbeitstag.

\subsection{SEPA-Lastschrift}

Bei der \emph{Lastschrift} zieht der Empfänger das Geld vom Konto des Zahlers ein ---
mit dessen Mandat (Einwilligung).

\begin{itemize}[itemsep=2pt]
  \item \textbf{SEPA-Basis-Lastschrift} (Core): typisch B2C, Erstattungsrecht 8 Wochen ohne Begründung, 13 Monate bei nicht autorisierter Lastschrift.
  \item \textbf{SEPA-Firmen-Lastschrift} (B2B): zwischen Unternehmen, kein Erstattungsrecht --- höhere Sicherheit für den Empfänger.
\end{itemize}

\subsection{Dauerauftrag}

Wiederkehrende Überweisung mit gleichem Betrag und Empfänger (z.\,B.\ Miete). \emph{Vom
Zahler initiiert} --- im Gegensatz zur Lastschrift.

\section{Zahlungskarten}

\begin{center}
\begin{tabular}{p{3.5cm} p{10.5cm}}
\toprule
\textbf{Karte} & \textbf{Funktion}\\
\midrule
Debitkarte (girocard) & Sofortige Belastung des Kontos --- Kontodeckung erforderlich.\\
Kreditkarte & Verfügungsrahmen, monatliche Sammelabrechnung; Charge (volle Zahlung) oder Revolving (Teilzahlung mit Zinsen).\\
Prepaid-Karte & Vorab aufgeladenes Guthaben --- ideal für Jugendliche oder Online-Shopping.\\
\bottomrule
\end{tabular}
\end{center}

\section{Onlinebanking, Mobil-Banking, TAN-Verfahren}

Onlinebanking nutzt \textbf{Zwei-Faktor-Authentifizierung} (PSD2): Ein zweites,
unabhängiges Sicherheitsmerkmal neben PIN/Passwort.

\begin{itemize}[itemsep=2pt]
  \item \textbf{chipTAN}: TAN-Generator mit Bankkarte erzeugt eine einmal gültige TAN.
  \item \textbf{photoTAN}: Farbiger Code im Banking-Programm wird mit App fotografiert.
  \item \textbf{pushTAN}: TAN per App (eigene Banking-App, durch separate App freigegeben).
  \item \textbf{Biometrie}: Fingerabdruck oder Face-ID auf dem Smartphone.
\end{itemize}

\begin{merksatz}[title={Sicherheitsregeln}]
\begin{enumerate}[itemsep=0pt]
  \item Niemals TAN am Telefon weitergeben --- Banken fragen nie danach.
  \item Vor jeder TAN-Eingabe Empfänger und Betrag im Generator/App prüfen.
  \item Nur HTTPS-Verbindungen nutzen (Schloss-Symbol im Browser).
\end{enumerate}
\end{merksatz}

\section{Zahlungsformen im E-Commerce}

\begin{center}
\begin{tabular}{p{3cm} p{6cm} p{4.5cm}}
\toprule
\textbf{Form} & \textbf{Beschreibung} & \textbf{Bewertung}\\
\midrule
Vorkasse & Zahlung vor Versand & Sicher für Händler, riskant für Käufer\\
Rechnung & Zahlung nach Lieferung & Käuferfreundlich, hohes Ausfallrisiko Händler\\
Lastschrift & Einzug nach Lieferung & Praktisch, Mandat erforderlich\\
Kreditkarte & Sofortige Autorisierung & Schnell, weltweit, Käuferschutz oft enthalten\\
PayPal & Wallet-Zahlung & Käuferschutz, Gebühren für Händler\\
Klarna & Rechnungs- oder Ratenkauf via Klarna & Händler bekommt Geld sofort, Käufer bezahlt später\\
Sofortüberweisung & Echtzeit-Überweisung über Klarna-System & Schnell, aber Bankzugangsdaten teilen\\
\bottomrule
\end{tabular}
\end{center}

\section{Skonto rechnen --- der IHK-Klassiker}

\subsection{Was ist Skonto?}

\emph{Skonto} ist ein Preisnachlass, den der Lieferant gewährt, wenn der Käufer
innerhalb einer kurzen Frist zahlt. Er belohnt schnelle Zahlung und reduziert das
Forderungsausfallrisiko des Lieferanten.

\subsection{Beispiel-Zahlungsbedingung}

\textit{,,Zahlbar in 30~Tagen netto Kasse oder innerhalb 14~Tagen mit 2\,\% Skonto.''}

\begin{beispiel}
Rechnung netto 5\,000~€. Skonto-Inanspruchnahme:
\begin{itemize}[itemsep=0pt]
  \item Skontobetrag $= 5\,000 \times 0{,}02 = \mathbf{100{,}00~€}$
  \item Zahlbetrag $= 5\,000 - 100 = \mathbf{4\,900{,}00~€}$
\end{itemize}
\end{beispiel}

\begin{merksatz}[title={Effektivverzinsung der Skonto-Inanspruchnahme}]
\textbf{Effektivverzinsung} $= \dfrac{\text{Skontosatz} \times 360}{\text{Zahlungsfrist} - \text{Skontofrist}}$
\end{merksatz}

\begin{beispiel}
Bei 2\,\% Skonto, 14~Tage Skontofrist, 30~Tage Zahlungsziel:
$\dfrac{2 \times 360}{30 - 14} = \dfrac{720}{16} = \mathbf{45{,}00\,\%~\text{p.\,a.}}$

Selbst wenn der Bankkredit 9~\% kostet --- Skonto ist mit 45~\% effektiv ungleich
teurer abzulehnen. Praxisregel: \textbf{Skonto fast immer ziehen.}
\end{beispiel}

\subsection{Was ist die Skonto-Bezugsgröße?}

Skonto wird typischerweise vom \textbf{Rechnungsbetrag nach Abzug von Liefererrabatten}
und ggf.\ vor USt berechnet. In der Klausur wird der Bezug klar genannt --- ,,vom
Rechnungsbetrag'', ,,vom Nettowarenwert'' usw.

\subsection{,,2/10 --- netto 30''}

Eine in den USA verbreitete, auch in Deutschland gebräuchliche Kurzschreibweise:
\textbf{2~\% Skonto bei Zahlung innerhalb 10~Tagen, sonst voller Betrag fällig in 30~Tagen.}

Effektivverzinsung: $\dfrac{2 \times 360}{30 - 10} = 36\,\%$ p.\,a.

\section{Zusammenfassung}

\begin{enumerate}[itemsep=2pt]
  \item \textbf{Geld}: Funktionen Tausch / Wertaufbewahrung / Maßstab; Bargeld vs. Buchgeld.
  \item \textbf{Zahlungsarten}: bar, halbbar, bargeldlos.
  \item \textbf{SEPA}: Überweisung, Basis-/Firmen-Lastschrift, Dauerauftrag --- IBAN-basiert.
  \item \textbf{Karten}: Debit, Kredit (Charge/Revolving), Prepaid.
  \item \textbf{Onlinebanking}: chipTAN, photoTAN, pushTAN; immer Zwei-Faktor.
  \item \textbf{E-Commerce}: Vorkasse, Rechnung, Lastschrift, Karte, PayPal, Klarna, Sofortüberweisung.
  \item \textbf{Skonto}: Effektivverzinsung $= \frac{\text{Skontosatz} \times 360}{\text{Zahlungsfrist}-\text{Skontofrist}}$ --- meist deutlich über Bankzinsen.
\end{enumerate}

\newpage

\section*{Glossar}
\addcontentsline{toc}{section}{Glossar}

\glossareintrag{Bargeld}{Münzen und Banknoten --- gesetzliches Zahlungsmittel.}

\glossareintrag{Buchgeld (Giralgeld)}{Guthaben auf Bankkonten --- per Überweisung, Lastschrift, Karte verfügbar.}

\glossareintrag{Halbbare Zahlung}{Mischform: bar einzahlen, bargeldlos überweisen (z.\,B.\ Bareinzahlung am Bankschalter).}

\glossareintrag{SEPA}{Single Euro Payments Area --- einheitlicher Euro-Zahlungsraum (36 Länder).}

\glossareintrag{IBAN}{International Bank Account Number --- in Deutschland 22-stellig (z.\,B.\ DE89\ldots).}

\glossareintrag{BIC}{Bank Identifier Code --- internationale Bankleitzahl (heute meist nicht mehr nötig in der EU).}

\glossareintrag{SEPA-Basis-Lastschrift}{Lastschriftverfahren für B2C --- 8 Wochen Erstattungsrecht ohne Begründung.}

\glossareintrag{SEPA-Firmen-Lastschrift}{Lastschriftverfahren zwischen Unternehmen --- kein Erstattungsrecht.}

\glossareintrag{Dauerauftrag}{Wiederkehrende Überweisung mit festem Betrag und Empfänger --- vom Zahler initiiert.}

\glossareintrag{Debitkarte (girocard)}{Karte mit sofortiger Kontobelastung --- nur bei Deckung.}

\glossareintrag{Kreditkarte}{Karte mit monatlicher Sammelabrechnung; Charge (volle Zahlung) oder Revolving (Teilzahlung).}

\glossareintrag{Prepaid-Karte}{Karte mit vorab aufgeladenem Guthaben.}

\glossareintrag{TAN}{Transaktionsnummer --- Einmal-Code zur Freigabe einer Bankaktion.}

\glossareintrag{chipTAN}{TAN-Verfahren mit kartenbasiertem Generator (Lesen eines Flickercodes vom Bildschirm).}

\glossareintrag{photoTAN}{TAN-Verfahren mit Mobile-App, fotografiert einen farbigen Code im Banking-Browser.}

\glossareintrag{pushTAN}{TAN-Verfahren mit eigener App auf dem Smartphone.}

\glossareintrag{PSD2}{EU-Richtlinie zur Zahlungsdienste --- erfordert Zwei-Faktor-Authentifizierung.}

\glossareintrag{Vorkasse}{Bezahlung vor Lieferung --- sicher für Händler, riskant für Käufer.}

\glossareintrag{Kauf auf Rechnung}{Zahlung nach Lieferung --- käuferfreundlich, höheres Ausfallrisiko für Händler.}

\glossareintrag{PayPal}{Wallet-System mit Käuferschutz; Händler zahlen Transaktionsgebühr.}

\glossareintrag{Klarna}{Zahlungsdienstleister: Rechnungs- oder Ratenkauf für den Käufer, Händler erhält Geld sofort.}

\glossareintrag{Sofortüberweisung}{Echtzeit-Überweisung über Klarna --- Käufer gibt Online-Banking-Daten ein.}

\glossareintrag{Skonto}{Preisnachlass des Lieferanten bei schneller Zahlung (typisch 2--3\,\% bei 10--14 Tagen Frist).}

\glossareintrag{Skontofrist}{Frist, innerhalb der gezahlt werden muss, um Skonto zu erhalten (z.\,B.\ 14 Tage).}

\glossareintrag{Zahlungsziel}{Maximale Zahlungsfrist ohne Skonto (z.\,B.\ 30 Tage netto Kasse).}

\glossareintrag{Effektivverzinsung der Skonto-Inanspruchnahme}{$=\frac{\text{Skontosatz} \times 360}{\text{Zahlungsfrist}-\text{Skontofrist}}$ --- meist sehr hoch (30--60\,\%).}

\end{document}
